\renewcommand{\Anno}{2006}

\renewcommand{\NumeroProblema}{3}

\renewcommand{\Testo}{È ben noto che più ci si alza rispetto alla
	superficie terrestre e più è possibile vedere (o essere visti da) lontano.
	\begin{enumerate}
		\item Un satellite si trova ad altezza $h$ sulla verticale di un punto $P$ della superficie terreste.  Determinare la massima distanza da $P$ sulla superficie terrestre alla quale è ancora possibile vedere il satellite.
	
		\item Una ditta di telecomunicazioni vorrebbe piazzare 4 satelliti in modo che da ogni punto della superficie terrestre sia possibile vederne almeno uno.  Per ragioni tecniche, la massima altezza da terra alla quale i satelliti possono essere lanciati è uguale al doppio del raggio terrestre.
		
		Determinare se la ditta può raggiungere il suo scopo.
	
		\item Determinare se, in assenza del vincolo sulla massima altezza raggiungibile, sarebbe possibile disporre opportunamente 3 satelliti in modo che da ogni punto della superficie terrestre sia possibile vederne almeno uno.
	\end{enumerate} 
	
	{\small Nota: si approssimi la Terra con una sfera di raggio $r$.}}

\renewcommand{\Risposta}{\begin{enumerate}
	\renewcommand{\labelenumi}{(\alph{enumi})} 
	\item[(a)]  La massima distanza è
	$$r\cdot\arccos\frac{r}{r+h}.$$

	\item[(b)]  La ditta può raggiungere il suo scopo.

	\item[(c)]  Non è comunque possibile.
\end{enumerate}
}

\renewcommand{\Aiutino}{\begin{enumerate}
	\renewcommand{\labelenumi}{(\alph{enumi})} 
	\item[(a)]  Considerare la seguente figura
	\begin{center}
		\psset{unit=7ex}
		\pspicture(-4,-1)(1,1)
		\pstGeonode[PosAngle={-45,135,135}](0,0){O}(-1,0){P}(-2.5,0){S}
		\pstHomO[HomCoef=0.5,PointSymbol=none]{S}{O}{C}
		\pstInterCC[PosAngle=90,PointSymbolB=none]{C}{S}{O}{P}{H}{H'}
		\pstCircleOA[linecolor=blue]{O}{P}
		\pcline(S)(P)\bput{:U}{$h$}
		\pcline(P)(O)\bput{:U}{$r$}
		\pcline(H)(O)\aput{U}{$r$}
		\pcline(S)(H)
		\pstRightAngle{S}{H}{O}
		\endpspicture

	\end{center}
	in cui la massima distanza richiesta è la lunghezza dell'arco $PH$.

	\item[(b)]  Disporre i 4 satelliti ai vertici di un tetraedro regolare di cui la terra sia la sfera inscritta. Quanto misura il raggio della sfera inscritta ad un tetraedro regolare?

	\item[(c)]  Tre satelliti determinano un piano. Da quale(i) punto(i) della sfera questo piano non è visibile?
\end{enumerate}
}

\renewcommand{\SoluzioneUfficiale}{\DomandaA Consideriamo la seguente figura
	\begin{center}
		\psset{unit=7ex}
		\pspicture(-4,-1)(1,1)
		\pstGeonode[PosAngle={-45,135,135}](0,0){O}(-1,0){P}(-2.5,0){S}
		\pstHomO[HomCoef=0.5,PointSymbol=none]{S}{O}{C}
		\pstInterCC[PosAngle=90,PointSymbolB=none]{C}{S}{O}{P}{H}{H'}
		\pstCircleOA[linecolor=blue]{O}{P}
		\pcline(S)(P)\bput{:U}{$h$}
		\pcline(P)(O)\bput{:U}{$r$}
		\pcline(H)(O)\aput{U}{$r$}
		\pcline(S)(H)
		\pstRightAngle{S}{H}{O}
		\pstMarkAngle[MarkAngleRadius=0.2,LabelSep=0.4]{H}{O}{S}{$\alpha$}
		\endpspicture

	\end{center}
in cui il punto $S$ rappresenta la posizione del satellite. I punti della superficie terrestre dai quali si vede il satellite sono quelli ottenuti ruotando l'arco $PH$ intorno all'asse $SO$. Di conseguenza la massima distanza da $P$ alla quale il satellite è visibile corrisponde alla lunghezza dell'arco $PH$, e cioè $r\alpha$. Per determinare l'angolo $\alpha$, ci basta osservare che nel triangolo rettangolo $OHS$ si ha che
$$\cos\alpha=\frac{OH}{OS}=\frac{r}{r+h}.$$

Pertanto la lunghezza dell'arco $PH$ è data da
$$r\cdot\arccos\frac{r}{r+h}.$$

Si noti che, come ci si poteva facilmente aspettare, tale lunghezza è nulla per $h=0$ e tende a $r\pi/2$ quando $h\to+\infty$.

\DomandaB La ditta può realizzare il suo scopo disponendo i 4
satelliti in modo che siano i vertici di un tetraedro regolare, di cui la terra sia la sfera inscritta.  In questo modo è chiaro che da ogni punto della superficie terrestre si vede almeno un satellite. Verifichiamo che in tale disposizione la distanza di ogni satellite dalla superficie terrestre è uguale esattamente al doppio del raggio terrestre.  Per far questo, ci basta dimostrare che la distanza dei vertici di un tetraedro dal centro della sfera inscritta è uguale al triplo del raggio della sfera stessa, il che a sua volta è equivalente a dimostrare che il raggio della sfera inscritta ad un tetraedro è uguale ad 1/4 dell'altezza.  Questo è un fatto ben noto.  Per dimostrarlo, indichiamo con $h$ l'altezza, con $r$ il raggio della sfera inscritta, e con $T$ l'area di ciascuna faccia del tetraedro.  Il volume del tetraedro è allora $Th/6$.  D'altra parte, possiamo scomporre il tetraedro come unione di 4 tetraedri, ciascuno avente una faccia in comune con il tetraedro originario, ed il restante vertice nel centro della sfera.  Ciascuno di questi tetraedri più piccoli ha volume $Tr/6$, e pertanto
$$\frac{Th}{6}=4\cdot\frac{Tr}{6},$$
da cui $h=4r$.

\DomandaC Mostriamo che con 3 satelliti, comunque disposti, c'è
sempre almeno un punto della superficie terrestre dal quale non se ne vede nessuno.  Consideriamo infatti il piano che contiene i 3 satelliti (se sono addirittura allineati, consideriamo un qualunque piano che li contenga).  Consideriamo ora la retta passante per il centro della terra e perpendicolare al piano: tale retta intersecherà la superficie terrestre in 2 punti $A$ e $B$.  Se il piano interseca propriamente (non tangente) la superficie terrestre, allora da $A$ e da $B$ non si vede nessun punto del piano, dunque in particolare non si vedono i satelliti (figura di sinistra). Se il piano non interseca la superficie terrestre, o è ad essa tangente, allora il piano è invisibile dal punto tra $A$ e $B$ che è più lontano dal piano stesso (figura di destra).

\psset{unit=5ex}
\pspicture(-4,-2)(4,2)
\pscircle(0,0){1}
\psline[linecolor=blue](-4,-0.3)(4,-0.3)
\pstGeonode[PosAngle={90,-90}](0,1){A}(0,-1){B}
\psline[linecolor=green](A)(-3.5,-0.3)
\psline[linecolor=green](A)(-2.5,-0.3)
\psline[linecolor=green](A)(-1.5,-0.3)
\psline[linecolor=green](B)(-3,-0.3)
\psline[linecolor=green](B)(-2,-0.3)
\endpspicture
\hfill
\psset{unit=5ex}
\pspicture(-4,-2)(4,2)
\pscircle(0,0){1}
\psline[linecolor=blue](-4,-1.5)(4,-1.5)
\pstGeonode[PosAngle={90,90}](0,1){A}(0,-1){B}
\psline[linecolor=green](A)(-3.5,-1.5)
\psline[linecolor=green](A)(-2.5,-1.5)
\psline[linecolor=green](A)(-1.5,-1.5)
\endpspicture

}

\renewcommand{\SoluzioniAlternative}{
}

\renewcommand{\Osservazioni}{\TitoloOsservazione }

\renewcommand{\Commenti}{
}